% Pursuing Stacks (cotes 134-2 to 134-8): archive cote 134-7, site page 70 (PDF page 71), typescript page 512.
% typescript. Draft compared: ps6.tex:3468-3545.
% AI-transcribed by Claude Opus 5.5 (claude-opus-5-5, adaptive thinking, effort high), Message Batches API, batch ps4 (msgbatch_012Aoc1NuskATPsyBeKSYK7J), 2026-11-12; not yet checked by eye.
\scanpage{134-7/70}
\mspage{512}

Moreover the canonical map
\[ X \hand{\longrightarrow} W(X) = \underline{Z}^{(X)} \]
factors (by construction) through
\[ (3) \qquad X \hand{\longrightarrow} W(X)(1) \quad ( \hand{\hookrightarrow} W(X) ) , \]
and this map is universal for all maps of X into torsors under abelian
group objects of $\hand{\Delta}^{\hand{\wedge}}$ . It is an immediate consequence of Whiteheads
theorem tha\over{r}{t} a weak equivalence $X \hand{\to} Y$ in $\hand{\Delta}^{\hand{\wedge}}$ induces \xout{\unsure{w}} weak equivalen-%
ce\hand{s} in $\hand{\Delta}^{\hand{\wedge}}$
\[ (4) \qquad \overset{\hand{\sim}}{W}(X) \hand{\longrightarrow} \overset{\hand{\sim}}{W}(Y) \quad , \qquad W(X)(1) \hand{\longrightarrow} W(Y)(1) \]
(and also $W(X)(n) \hand{\to} W(Y)(n)$ for any $n \hand{\in} \underline{Z}$ ), and accordingly that
the map between the normalized chain complexes corresponding to $\overset{\hand{\sim}}{W}(X)$ and
$\overset{\hand{\sim}}{W}(Y)$, namely by definition
\[ (4') \qquad L\overset{\hand{\sim}}{H}.(X) \hand{\longrightarrow} L\overset{\hand{\sim}}{H}.(Y) \]
is a quasi-isomorphism. Thus we get \xout{functors} a functor
\[ \overset{\hand{\sim}}{\mathrm{LH}}. : \mathrm{Hot} \hand{\longrightarrow} \mathrm{Hotab} \quad , \]
whose composition\inserted{$L\widetilde{H}.\circ\beta$ ,} with the forgetful functor
\[ \hand{\beta} : \mathrm{Hot}^{\hand{\bullet}} \hand{\longrightarrow} \mathrm{Hot} \]
is \xout{the funct} canonically isomorphic to the functor $L\overset{\hand{\sim}}{H}.$ of p.510 (formula
9). More specifically, when X is in $\hand{\Delta}^{\hand{\wedge}\hand{\bullet}}$ , then the map
\[ \underline{Z} = \underline{Z}^{(e)} \hand{\longrightarrow} \underline{Z}^{(X)} \quad \text{deduced from } e \hand{\to} X \]
defines a splitting of the extension 1, i.e. an isomorphism
\[ (5) \qquad \underline{Z}^{(X)} = W(X) \hand{\simeq} \overset{\hand{\sim}}{W}(X) \hand{\oplus} \underline{Z} \quad , \]
and accordingly $\overset{\hand{\sim}}{W}(X)$ may be interpreted, at will, as a quotient group
of W(X) (which we did on p.510)\hand{,} or as a subobject of W(X),\deleted{(}which is the
better choice, because it makes a good sense without using any pointed
structure. Again, without using the pointed structure, we\xout{r} get a canoni-%
cal exact trian\xout{\unsure{g}}gle in D(Ab) interpreting (1)
% Hand-drawn arrows are marked blue; typed objects remain plain.
\diagram{
\begin{tikzcd}[column sep=small]
 & \underline{Z} \arrow[color=blue,dl] & \\
L\overset{\hand{\sim}}{H}.(X) \arrow[color=blue,rr] & & LH.(X) \arrow[color=blue,ul]
\end{tikzcd}
}
(6)

replacing ((12) p. 510), which \deleted{should be replaced} (in the case considered
there, namely X pointed\hand{)} \inserted{should be replaced} by the more precise relationship
\[ (\xout{4}\, 5') \qquad LH.(X) \hand{\simeq} L\overset{\hand{\sim}}{H}.(X) \hand{\oplus} \underline{Z} \quad , \]
an isomorphism functorial for X in Hot\hand{$^{\bullet}$}.

\margin{70}

% Visual collation: both supplied image halves inspected; see CHECK_DECISIONS_134_7_8.tsv.
